\section{The paper and the formal theorem: correspondence and differences}
\label{sec:paper-correspondence}

The principal result formalized here is the general sufficiency-of-robustness
argument. Its reference in the
\href{https://arxiv.org/pdf/1907.01257v5\#page=30}{version~5 PDF of the paper}
is \emph{Theorem~8.11}. The
\href{https://arxiv.org/html/1907.01257v5\#S8.SS4}{HTML rendering}
numbers the same result Theorem~6. The Lean name \texttt{Unbounded.theorem6}
and many references in this commentary follow the HTML numbering. They do not
identify a different mathematical result. This section uses the PDF numbering
when discussing the original statement.

The Lean theorem proves the two corresponding contextual-refinement
implications for the definitions constructed in this repository. It is not a
literal transcription of the paper. Some changes are representation choices
with proved comparison lemmas; others require a further audit of the
correspondence with the paper. In particular, the missing plugging-closure
premise and the treatment of substitution stability, discussed below, should
not be silently described as harmless changes or established strengthenings.
A kernel-checked proof establishes the formal proposition; it does not by
itself establish that the proposition expresses every condition intended in
the paper.

\Needspace{14\baselineskip}
\subsection{The statement and its quantitative conclusion}

Write $T(G,H)$ for the pre-template relation and $G\preceq_Q^{\mathcal C}H$
for contextual refinement with cost comparison $Q$. The forward conclusion
has the form
\[
 T(G,H)\quad\Longrightarrow\quad G\preceq_Q^{\mathcal C}H,
\]
under input safety and robustness. The second implication reverses $G,H$;
it uses input safety with the converse cost relation and the corresponding
converse robustness condition. Both clauses share output closure,
deterministic and refocusing behaviour, and a reasonable triple of cost
preorders. In the final Lean declaration the paper's $(Q,Q',Q'')$ are named
$(Q,P,R)$. These are parameters, not specific numerical inequalities.
\lean{Graph/UnboundedTheorem}{110}{124}{Unbounded.theorem6}

State refinement compares successful termination and exact numbers of machine
steps: a terminating execution on the source must have a terminating matching
execution on the target whose lengths satisfy the chosen preorder. The
conclusion does not assert equality of final graphs or values. Contextual
refinement makes this comparison after observation contexts are filled.
Consequently the meaning of the conclusion depends on the definitions of
admissible observations, as well as on the transition relation.

\Needspace{14\baselineskip}
\subsection{Hypernets, labels, and equality}

The paper presents finite typed hypergraphs with interfaces and recursively
boxed graphs. Lean separates internal vertices from indexed input and output
ports, makes the ordered interface types explicit, and carries finiteness and
incidence laws in the graph structure. This replaces informal choices of
names and boundary identifications with functions, finite indices, and proof
fields. Plugging therefore requires explicit interface transport, even when
its graphical meaning is immediate.

A finite allowance for nesting makes the recursive construction of boxed
labels well founded. The allowance is not the intrinsic depth of the graph:
the same graph can be represented at larger allowances. Global graph classes
identify recursive graph isomorphism and such padding. Recursive isomorphism
also permits independent renaming inside boxes. The relevant equality is thus
structural equality up to these identifications, rather than literal equality
of the chosen vertex and edge carriers.

There is a further, substantive label issue. The internal typed representation
records box-boundary information that cannot simply be erased: the paper
requires equal labels to have equal signatures. The development gives a
freshening encoding into actual graph-valued labels, with encoding/decoding
round trips on recursive graph classes. These are proved representation
results; the earlier counterexample to naive erasure explains why a direct
forgetful translation would not have sufficed.
See \leanfile{Hypernets/Graph/PaperEncoding.lean}{PaperEncoding.lean} and
\leanfile{Hypernets/Graph/PaperRelations.lean}{PaperRelations.lean}.

\Needspace{14\baselineskip}
\subsection{States and genuinely unbounded execution}

A state is a suitably focussed closed graph. Focus insertion and erasure let
the development reason using a focus-free graph and a distinguished coordinate,
then reconstruct the actual focussed graph. Search is first expressed in
these coordinates. Separate adequacy proofs compare it in both directions
with adjacency equations and with the finite interaction fragments in simple
hole contexts. Coordinates are a proof device, not an additional machine
operation. The word \emph{corolla} denotes retained-edge bookkeeping used in
this construction; it introduces no additional semantic rule from the paper.
See \leanfile{Hypernets/Graph/HoleContextSearch.lean}{HoleContextSearch.lean}.

The earlier fixed-allowance machines were intermediate results. The final
machine uses the union of all finite allowances, quotiented by renaming and
padding. A finite run can be represented at a sufficiently large common
allowance with exactly the same length. No single allowance bounds all runs,
and computation may increase the necessary allowance. Finality and the
full transition relation are transferred, not just search.
See \leanfile{Hypernets/Graph/UnboundedUAM.lean}{UnboundedUAM.lean}.

Rootedness means reachability by search from an initial state; it does not
mean arbitrary reachability by the full machine. The global rooted predicate
is proved equivalent to this search-reachability formulation. Refocusing of
extrinsic behaviour means preservation of rootedness by its compute steps;
the full-machine preservation theorem combines this assumption with the
intrinsic search and substitution results.
See \leanfile{Hypernets/Graph/UnboundedRoots.lean}{UnboundedRoots.lean}.

The parameter \texttt{Behaviour} supplies the extrinsic compute relation.
Determinism and refocusing of that behaviour remain assumptions of the general
theorem, as do the semantic conditions on a proposed template. The development
proves the required structural results for the assembled machine. Equality of
successors in its determinism theorem is equality of state classes, expressing
determinism modulo the chosen structural identifications.

\Needspace{14\baselineskip}
\subsection{Templates, specimens, and robustness}

A global pre-template is a relation between focus-free graph classes with
matching interfaces. Its interpretation at every finite allowance gives the
layered pre-template used by the concrete constructions. What the paper groups
as template conditions is supplied through separate Lean premises: in
particular, output closure and input safety are not consequences of an
arbitrary pre-template relation.

A specimen consists of an actual context together with two related fillings.
The two resulting states share the distinguished context focus. The formal
representation uses named, typed slots and recursively fills the original
holes. Inserted fragments are not recursively substituted again during the
same filling operation. Although the slot environment can declare unused
slots, a finite graph has only finitely many actual hole occurrences. Tuple
representation lemmas connect these fillings with finite families of original
holes; the slot environment itself should not be confused with an infinite
graph or an infinite simultaneous rewrite.
See \leanfile{Hypernets/Graph/ObservationFillings.lean}{ObservationFillings.lean}.

Quasi-specimens retain an unconditional specimen witness and agreement of
focus modes. Only the rootedness and comparison obligations are conditional
on the two endpoint states being rooted. Input safety retains its three
alternatives. Robustness likewise allows an unrooted case, a stuck continuation,
or a matching continuation through a non-rewrite quasi-specimen, with the
specified costs and up-to comparisons. In the final definitions, stuckness
means stuckness in the unbounded machine. It is not inferred merely from
being stuck at one finite allowance.
\lean{Graph/UnboundedSafety}{25}{59}{Safety.Quasi, InputSafe, Robust}

\Needspace{14\baselineskip}
\subsection{Observations and the context-closure discrepancy}
\label{sec:context-correspondence-audit}

The formal predicate \texttt{PaperRefines} quantifies over every pair of finite
representatives of the two global graphs, every hole name, and the observation
contexts specified independently of the template. The graph interfaces must
match, and the observation construction imposes its own well-formedness and
closure requirements. A proved common-representatives lemma ensures that the
universal quantification over representatives is not vacuous. This lemma does
not, on its own, prove that every observation intended by the paper is covered.
\lean{Graph/UnboundedTheorem}{9}{17}{PaperRefines}
\lean{Graph/UnboundedTheorem}{30}{41}{common_representatives}

There is an explicit difference in hypotheses. The paper's contextual setup
in Section~7.1, and Theorem~8.11 itself, require the context set $\mathcal C$
to be closed under plugging: whenever the relevant composite of allowed
contexts is defined, it is again allowed. The final Lean theorem instead
takes an arbitrary predicate $C$ on global context classes. It induces
membership at each finite allowance by interpreting the finite graph in the
global class. The resulting membership is invariant under graph isomorphism
and padding.
\lean{Graph/UnboundedTheorem}{89}{105}{globalAllowed, globalAllowed_iso, globalAllowed_padding}

\emph{Invariance under renaming or padding is not closure under plugging.}
The latter is not a premise of \texttt{theorem6}. Its absence does not invalidate
the proved Lean implication, but it prevents treating the two statements as
identical without further analysis. A stronger theorem for arbitrary context
sets would imply the instance for plugging-closed sets \emph{if} the observation
predicates agree. The remaining audit is to establish that agreement and
explain exactly where the paper's closure condition is used, or to adjust the
formal setup if its observation quantification differs. This commentary does
not claim that this audit has already been completed.

\Needspace{14\baselineskip}
\subsection{Substitution and the stability discrepancy}
\label{sec:stability-correspondence-audit}

The paper's Lemma~7.6 uses stability of the substituted graph in its argument
reducing full-machine refocusing to the extrinsic behaviour. The Lean
substitution-rootedness theorem has no explicit stability premise for that
graph. Its proof instead verifies a stationary-plugging criterion using the
actual variable-prefix geometry: the variable prefix blocks an operation path
from the sole input, and the condition concerning other inputs is vacuous
for a unary interface.
\lean{Graph/SubstitutionRefocusing}{25}{39}{Substitution.prefix_oneWay}
\lean{Graph/SubstitutionRefocusing}{76}{104}{State.substitution_rooted}

This is a proved result about the encoded substitution rule. It is not yet a
justification for dropping the paper's stability condition. The correspondence
audit must compare the admissible instances of Figure~9(f), the definition of
rootedness, and the stationary-plugging criterion with the paper's argument.
Only after that comparison can the difference be classified as a valid
strengthening, a condition supplied elsewhere, or a modelling discrepancy.
The final theorem's use of the formal refocusing result makes this comparison
material to its relation with Theorem~8.11.

\Needspace{14\baselineskip}
\subsection{Proof organization and the boundary of the claim}

The proof retains the paper's main route: local specimen analysis yields a
counting simulation; soundness of that simulation yields state refinement;
initial observation lifting yields contextual refinement. Lean factors the
first part through abstract machine and specimen-system interfaces, then
constructs concrete instances and proves their laws. The abstraction does not
leave those concrete laws as unproved axioms. Additional transport, quotient,
freshness, and common-allowance lemmas account for much of the formal proof's
length; they make representation choices explicit rather than replace the
central simulation argument.

The kernel checks the final propositions and their dependencies. The recorded
axiom audits exclude project-specific axioms and unfinished proofs; use of
Lean's standard classical and quotient principles is part of the foundation.
None of these checks automatically validates the translation of informal
mathematics into definitions.

Accordingly, the completed milestone is the unbounded contextual sufficiency
theorem \emph{in the repository's concrete encoding}. The two correspondence
questions above remain distinct from missing Lean proof terms and are recorded
in Appendix~\ref{app:remaining-work}. The copying UAM, translations of specific
languages, and application-specific equivalences were outside that general-theorem
milestone. The new target is Proposition~12.2: Section~\ref{sec:copying-foundations}
begins its separate copying/store development and states what remains.
Sections~\ref{sec:copying-assembly}--\ref{sec:copying-values} construct the
distributors, tree/store tracing, full Figure~16 translation, and stable/copyable
graph predicates, preserving the proposition's terms and all prescribed
distributors. These structural results do not yet supply the copying machine
or the observational equivalence of Proposition~12.2.
Section~\ref{sec:copying-binding} proves sourced-hole binding-free plugging
closure and exhibits a zero-input-hole counterexample to unrestricted closure
under the literal Definition~11.2. Translated term holes always have an input;
this qualification does not change the proposition's terms or their translation.
Section~\ref{sec:copying-observations} adds annotated term contexts and a
concrete closing-context observation interface parameterised by execution.
Exact constructor-interpretation laws are proved.
Section~\ref{sec:copying-context-binding} proves binding-freedom of all
translated lambda-contexts. Section~\ref{sec:copying-certified-contexts}
adds actual sole-hole addresses, composition of graph contexts, and
graph-context congruence conditional on execution invariance under recursive
isomorphism. Section~\ref{sec:copying-expansion-plug} proves actual
addressed graph-filling correspondence for unary term contexts.
Section~\ref{sec:copying-nested-filling} extends this to arbitrarily nested
lambda/unary contexts, with sole-hole and binding-free certificates.
Section~\ref{sec:copying-parallel-padding} proves filling through both parallel
arguments, the full application/binary clauses, and arbitrary iterated padding.
Section~\ref{sec:copying-context-filling} proves unrestricted addressed
term-context correspondence, including cumulative padding at common-depth
joins and a single certified address valid for all fitting graphs.
Section~\ref{sec:copying-observation-padding} proves observation reflection
from larger allowances and the cofinal comparison criterion, with explicit
execution-isomorphism and padding hypotheses. A deterministic counterexample
shows that the latter is an independent obligation.
Section~\ref{sec:copying-term-congruence} derives observational congruence for
all original contexts when the input terms have equal minimal depths, assuming
execution isomorphism invariance. The proposition's actual depths are three
and two. Section~\ref{sec:copying-general-congruence} removes the depth
restriction and proves general observational context congruence under the
explicit execution isomorphism and padding laws. Those concrete machine
laws remain open; no equivalence of the proposition programs is inferred.
Earlier statements that no work remained for the paper's theorem should be
read with this qualification: completion of the formal construction does not
settle these external correspondence questions.
\clearpage


Section~\ref{sec:copying-padding-predicates} proves that the actual stability
predicate is preserved and reflected by padding, and copyability is preserved.
These structural laws do not discharge the copying execution's representation
laws; they retain the distinction between stable graphs and copyable payloads.


Section~\ref{sec:copying-translated-composition} establishes actual graph
composition of independently translated contexts. The original hosts,
interfaces, term constructors and translations are retained; the correspondence
is recursive graph isomorphism at equal outer depth, without an execution premise.


Section~\ref{sec:copying-atomic-renaming} proves selected-resource renaming and
all atomic translation cases for finite resource permutations, retaining the
original weakening/contraction forests and instance prefix. These are actual
boundary isomorphisms of the original translations. They do not establish
composite-term renaming or the final observational equivalence.
